\chapter*{}
%\thispagestyle{empty}
%\cleardoublepage
%\thispagestyle{empty}

% \input{portada/portada_2}


%\cleardoublepage
\thispagestyle{empty}

\begin{center}
{\large\bfseries \myTitle}\\
\end{center}
\begin{center}
\myName\\
\end{center}

%\vspace{0.7cm}
\noindent{\textbf{Palabras clave}: calibración de cámara, modelado geométrico, metrología monocular, aprendizaje semi-supervisado, visión por computador.
}\\

\noindent{\textbf{Resumen}}\\
La metrología monocular constituye una línea de investigación de gran relevancia debido a su potencial
para obtener mediciones precisas a partir de una única imagen, es decir, una sola vista de la escena.
Es especialmente útil en situaciones donde no es posible o práctico capturar múltiples imágenes,
desde diferentes ángulos. El auge del aprendizaje profundo ha abierto nuevas oportunidades para abordar
este problema, permitiendo inferir información a partir de patrones visuales aprendidos.
Este Trabajo de Fin de Máster presenta una implementación y adaptación del método de aprendizaje semi-supervisado descrito por Zhu \emph{et al.}
en \emph{Single View Metrology in the Wild} (publicado en la \emph{European Conference on Computer Vision} en 2020), debido a la ausencia de una implementación pública disponible.
\smallskip

La implementación se realiza en el reciente \emph{framework} de detección de objetos \emph{Detectron2},
permitiendo aprovechar arquitecturas más recientes, mejores mecanismos de optimización y entrenamiento distribuido, así como un rendimiento computacional superior.
En relación a los datos, se utiliza un conjunto de datos obtenido de la red social Flickr
para el entrenamiento del módulo de calibración de cámara, en sustitución del
conjunto de datos original SUN360 que ya no está disponible.
Además, se utilizan múltiples fotogramas de cámaras de seguridad
como conjunto de validación experimental.
\smallskip

El rendimiento del sistema se evalúa frente al procedimiento tradicional basado en la calibración manual de
la cámara y la obtención de medidas de referencia de la escena para estimar la altura de los objetos.
Si bien los resultados presentados en este trabajo son de carácter preliminar,
se obtuvo un error absoluto medio de $9.8\pm6.9$ cm.
Aunque el método implementado no supera al método de referencia, que presenta un error absoluto medio de $6.7\pm5.6$ cm,
elimina la necesidad de realizar una calibración manual específica de la escena y permite automatizar el proceso de estimación.
Asimismo, se identifican diversas líneas de mejoras para futuras investigaciones.

\cleardoublepage


\thispagestyle{empty}


\begin{center}
{\large\bfseries \myTitleENG}\\
\end{center}
\begin{center}
\myName \\
\end{center}

%\vspace{0.7cm}
\noindent{\textbf{Keywords}: camera calibration, geometric modelling, monocular metrology, semi-supervised learning, computer vision.}\\

%\vspace{0.7cm}
\noindent{\textbf{Abstract}}\\
Monocular metrology constitutes a highly relevant research area due to its potential to obtain accurate measurements from a single image, i.e.,
a single view of a scene. It is particularly useful in situations where capturing multiple images from different
viewpoints is not feasible or practical. The rise of deep learning has opened new opportunities to address this problem
by enabling the inference of geometric information from learned visual patterns.
This Master's Thesis presents an implementation and adaptation of the semi-supervised learning method described by Zhu \emph{et al.} in
\emph{Single View Metrology in the Wild} (presented at the European Conference on Computer Vision in 2020), due to the lack of a public implementation available.
\smallskip

The method is implemented using the Detectron2 object detection framework,
enabling the use of more recent architectures, improved optimization mechanisms and distributed training capabilities,
as well as superior computational performance. Regarding the data, a training dataset collected from the Flickr social network
is employed for the camera calibration module, as a replacement for the original SUN360 dataset which is no longer available.
In addition, multiple frames from surveillance cameras are used as an experimental validation dataset.
\smallskip

The performance of the system is evaluated against the traditional procedure based on manual camera calibration and the
acquisition of reference measurements within the scene to estimate object heights.
Although the results presented in this work are preliminary, the proposed method achieved a
mean absolute error of $9.8\pm6.9$ cm.
Although the implemented method does not outperform the reference method,
which achieves a mean absolute error of $6.7\pm5.6$ cm, it eliminates the need for scene-specific manual calibration and enables the estimation process to be automated.
Furthermore, several potential paths for improvement are identified for future research.

\chapter*{}
\thispagestyle{empty}

\noindent\rule[-1ex]{\textwidth}{2pt}\\[4.5ex]

Yo, \textbf{\myName}, alumno de la titulación Máster en Ciencia de Datos e Ingeniería de Computadores
de la \textbf{\myFaculty},
con Pasaporte NX4L843F5, autorizo la ubicación de la siguiente copia de mi
Trabajo Fin de Máster en la biblioteca del centro para que pueda ser
consultada por las personas que lo deseen.

\vspace{6cm}

\noindent Fdo: \myName

\vspace{2cm}

\begin{flushright}
Granada a \today.
\end{flushright}


\chapter*{}
\thispagestyle{empty}

\noindent\rule[-1ex]{\textwidth}{2pt}\\[4.5ex]

D. \textbf{\myProf}, Profesor del \myDepartment de la Universidad de Granada.

\vspace{0.5cm}

D. \textbf{\myOtherProf}, Profesor del \myDepartment de la Universidad de Granada.


\vspace{0.5cm}

\textbf{Informan:}

\vspace{0.5cm}

Que el presente trabajo, titulado \textit{\textbf{\myTitle}}
ha sido realizado bajo su supervisión por \textbf{\myName}, y autorizamos la defensa de dicho trabajo ante el tribunal
que corresponda.

\vspace{0.5cm}

Y para que conste, expiden y firman el presente informe en Granada a \today.

\vspace{1cm}

\textbf{Los directores:}

\vspace{5cm}

\noindent \textbf{\myProf \ \ \ \ \ \myOtherProf}

\chapter*{Agradecimientos}

\thispagestyle{empty}
       \vspace{1cm}


Primeramente, me gustaría agradecer a mis tutores, Enrique Bermejo y Pablo Mesejo,
por darme la oportunidad de desarrollar este proyecto con ellos.
Agradezco la paciencia infinita y comprensión a la hora de resolver mis dudas.
Segundo, agradezco a la propia Universidad de Granada por haberme dado la oportunidad
de continuar mis estudios universitarios en tan distinguida casa de estudios.
\\
